\documentclass{gescons}

\genre {Parceria}
\author{Ana Claudia do Prado e~Magda Stapf Amancio}
\authorrole{Coordenadoras Gerais da Editares}
\title{Editares e~Uniescon Firmam Acordo Inédito de Cooperação Interinstitucional}

\begin{document}
    \makeentrevistatitle
    %\maketitle

    %\fullwidthimage{fields}{b}

    %\coverart{back/editorial}
    \coverart{articles/parceria/fundo-parceria}
    
    
\begin{center}
    \includegraphics[height=4cm]{images/Logo-Editares-com-Marca-Registrada}
    \hspace{3cm}
    \includegraphics[height=5cm]{images/Logo-UNIESCON-2048x1741} 
    
    %\includegraphics[height=10cm]{articles/parceria/imagens/parceria-editares-uniescon}
    %\includegraphics[height=10cm]{example-image}
\end{center}
    
    \begin{multicols}{2}

A produção conscienciológica ganha um novo impulso. Em julho de 2025, a~Editares e~a~Uniescon firmaram um \textbf{acordo inédito de cooperação interinstitucional.} O~objetivo é~compartilhar os pareceres gesconográficos de autorandos, contribuindo para a~qualificação de obras em fase de construção e~revisão, de acordo com os trâmites editoriais das duas instituições.

O intercâmbio dos pareceres será realizado somente mediante consentimento prévio e~expresso do autorando, garantindo transparência e~respeito ao processo autoral. Editares e~Uniescon também assumiram o~compromisso de \textbf{preservar a~confidencialidade dos arquivos} e~manter sigilo sobre o~material que estiver fora do fluxo autorizado.

Com a~iniciativa, Editares e~Uniescon \emph{consolidam um trabalho conjunto em prol da melhoria contínua da produção da gesconografia conscienciológica,} criando novas pontes de apoio para pesquisadores e~escritores.



    \end{multicols}
\end{document}
